\begin{textsample}{POS Dim 1 – 2019\_c4 – Score -3.00 – 2019\_c4/enade\_2019\_agro\_96d666036e.txt}  \label{ex:f1_pos_341}
CAPÍTULO 4 PERCEPÇÃO DA PROVA Nas análises feitas neste capítulo, constam as percepções dos concluintes da Área de Agronomia sobre a prova aplicada no Enade/2019.
Essas percepções foram mensuradas por meio de nove questões que avaliaram desde o grau de dificuldade da prova até o tempo gasto para concluí -la.
As percepções sobre a prova foram relacionadas ao desempenho dos estudantes, à Grande Região de funcionamento do curso, à Categoria Administrativa e ao tipo de Organização Acadêmica da IES.
O questionário de percepção da prova encontra -se ao final do Anexo VII, onde está a reprodução do exame.
O desempenho dos estudantes foi classificado em quatro quartos.
Para tanto, esse desempenho foi ordenado de forma ascendente.
O percentil 25, P25, também conhecido como primeiro quartil, é a nota de desempenho que deixa um quarto ( 25 \% ) dos valores observados abaixo e três quartos acima.
Na Figura 1, apresenta -se uma ilustração desse conceito.
O quarto inferior de desempenho é composto pelas notas abaixo do primeiro quartil.
Já o percentil 75, P75, também conhecido como terceiro quartil, é o valor para o qual há três quartos ( 75 \% ) dos dados abaixo e um quarto acima dele.
O quarto superior de desempenho é composto pelas notas iguais ou acima do terceiro quartil.
O percentil 50, P50, também conhecido como mediana, é o valor que divide as notas em dois conjuntos de igual tamanho.
O segundo quarto inclui valores entre o primeiro quartil ( P25 ) e a mediana.
O terceiro quarto contém os valores entre a mediana ( P50 ) e o terceiro quartil ( P75 ).
Vale ressaltar que percentis, quartis e medianas são pontos que, não obrigatoriamente, pertencem ao conjunto original de dados, ao passo que os quartos são subconjuntos dos dados originais.
A seguir, serão apresentados gráficos com resultados selecionados, relativos às nove questões avaliadas por grupos de estudantes.
Nas barras dos gráficos, apresenta -se o percentual de alunos que assinalaram uma das opções ou a soma das porcentagens daqueles que assinalaram duas ( ou três ) delas.
Por exemplo, para as questões 1 e 2, nos gráficos, é apresentada a porcentagem total de participantes que assinalaram as opções ( D ) Difícil e ( E ) Muito difícil.
Em cada barra, foram assinalados, também, os extremos do intervalo de confiança de 95 \% como linhas verticais unidas por uma linha horizontal na forma da letra H maiúscula.
O estimador de um parâmetro com um certo nível de confiança ( e.g. 95,0 \% ) deve conter o parâmetro no intervalo de confiança em 95 \% das vezes.
Na comparação entre os estimadores dos parâmetros de duas classes de uma dada categoria ( e.g.
Norte e Nordeste nas Grandes Regiões, de primeiro e último quarto dentro de desempenho, de IES Públicas e Privada, ou de Universidades e Faculdades ) associados aos seus respectivos intervalos de confiança, diz -se que não há diferença estatisticamente significativa entre os parâmetros das duas categorias quando há interseção entre os intervalos de confiança e que há diferença, se os intervalos de confiança são disjuntos ( para maiores informações vide Glossário ).
Nas Tabelas do Anexo II, são apresentados os valores absolutos e a distribuição percentual18 das alternativas válidas das nove questões, segundo o mesmo recorte de desempenho dos alunos e Grande Região de funcionamento do curso, de Categoria Administrativa e do tipo de Organização Acadêmica da IES. 4.1 GRAU DE DIFICULDADE DA PROVA 4.1.1 Componente de Formação Geral Ao avaliarem Qual o grau de dificuldade desta prova na parte de Formação Geral? ( Questão 1 ), 38,1 \% do grupo de inscritos e presentes optaram pelas alternativas Difícil ou Muito difícil.
Entretanto, para mais da metade dos estudantes ( 56,3 \% ), o Componente de Formação Geral da prova foi considerado com grau de dificuldade Médio ( Gráfico 4.1, Gráfico 4.2 e, no Anexo II, a Tabela II.1 ).
O percentual de estudantes que considerou a prova Difícil ou Muito difícil foi maior na região Sul, onde a proporção foi 43,4 \%, enquanto na região Nordeste, a proporção foi a menor, o correspondente a 32,3 \%.
No Gráfico 4.1, é possível observar que as diferenças entre a região Sul e todas as outras as regiões são estatisticamente significativas, assim como as diferenças entre a região Nordeste e as regiões Sudeste e Centro-Oeste.
Nas Grandes Regiões, a proporção de presentes à prova que consideraram o Componente de Formação Geral de grau de dificuldade Médio esteve entre 52,0 \% na região Sul e 62,1 \% na região Nordeste.
O percentual de alunos que considerou o componente de Formação Geral da prova Difícil ou Muito difícil foi decrescente em função dos quartos de desempenho : de 41,9 \%, no primeiro quarto, e de 34,9 \%, no quarto, o grupo de melhor desempenho na prova.
Nos quartos de desempenho intermediários, a proporção de alunos que consideraram a prova Difícil ou Muito difícil foi de 39,2 \%, no segundo quarto e de 36,5 \%, no terceiro quarto.
A diferença entre os quartos de desempenho extremos, o primeiro e o quarto, é estatisticamente significativa.
Para todos os quartos de desempenho, a alternativa modal para essa pergunta foi a Médio, variando entre 53,1 \% e 58,1 \% dos respondentes no primeiro e terceiro quartos de desempenho, respectivamente.
O percentual de estudantes que considerou o componente de Formação Geral da prova Difícil ou Muito difícil foi maior entre os estudantes de Centros Universitários, alcançando 48,3 \%, enquanto os estudantes de Universidades atingiram 33,6 \%.
No Gráfico 4.3, a única diferença entre as proporções que não é estatisticamente significativa é entre as Faculdades e os CEFET/IF.
Nas Organizações Acadêmicas, a proporção de presentes à prova que avaliaram o Componente de Formação Geral como de grau de dificuldade Médio esteve entre 48,0 \% nos Centros Universitários e 60,0 \% nas Universidades ( ver também a Tabela II.2, no Anexo II ).
O percentual de estudantes que considerou o componente de Formação Geral da prova Difícil ou Muito difícil foi maior entre os estudantes de cursos de IES Privadas ( 46,3 \% ), uma diferença estatisticamente significativa da proporção alcançada em IES Públicas ( 32,7 \% ).
Já a proporção de presentes à prova que consideraram o Componente de Formação Geral de grau de dificuldade Médio foi 60,4 \% para os de IES Públicas e 50,1 \% para os de IES Privadas ( ver Gráfico 4.4 ). 4.1.2 Componente de Conhecimento Específico Ao responderem à Questão 2 – Qual o grau de dificuldade desta prova na parte de Componente Específico? – 32,3 \% do grupo de estudantes classificaram -na como Difícil ou Muito difícil.
Além disso, o Componente modal de Conhecimento Específico da prova foi avaliado com grau de dificuldade Médio por 61,3 \% dos alunos ( Gráfico 4.5, Gráfico 4.6 e, no Anexo II, a Tabela II.3 ).
A análise das respostas dos estudantes quanto ao grau de dificuldade do Componente de Conhecimento Específico da prova, agregado por Grande Região, mostra que a diferença entre a maior e todas as outras proporções de alunos que a avaliaram como Difícil ou Muito difícil é estatisticamente significativa, sendo a maior, na região Sul ( 36,0 \% ).
As diferenças entre as regiões Norte, Nordeste, Sudeste e Centro-Oeste não são estatisticamente significativas.
Os alunos que classificaram o grau de dificuldade como Médio, no Componente de Conhecimento Específico, atingiram um percentual que variou de 57,6 \% a 65,3 \%, para as regiões Sul e Nordeste, respectivamente.
Considerando -se a avaliação da dificuldade das questões do Componente de Conhecimento Específico da prova, de acordo com o desempenho dos estudantes, observa -se que há diferença estatisticamente significativa dos resultados entre o último quarto e os demais quartos de desempenho, assim como há diferença entre o penúltimo quarto e os dois primeiros quartos.
Já a diferença entre os dois primeiros quartos de desempenho não é estatisticamente significativa.
A proporção dos que classificaram a parte específica como Difícil ou Muito difícil variou de 25,2 \% ( último quarto ) a 36,6 \% ( primeiro quarto ).
As demais proporções dessa reposta foram de 35,6 \% e 31,8 \%, nos segundo e terceiro quartos, respectivamente.
Já a proporção dos que responderam que o grau de dificuldade das questões do Componente de Conhecimento Específico da prova foi Médio variou de 59,1 \% no primeiro quarto a 65,3 \% no último.
Considerando -se a avaliação da dificuldade das questões do Componente de Conhecimento Específico da prova, o percentual de estudantes que as considerou Difícil ou Muito difícil foi maior entre os estudantes de Centros Universitários, o equivalente a 40,6 \%, enquanto os de Universidades alcançaram a menor proporção entre as Organizações Acadêmicas : 28,8 \%.
No Gráfico 4.7, é possível observar que a única diferença que não é estatisticamente significativa é entre as Faculdades ( 36,3 \% ) e os CEFET/IF ( 32,8 \% ).
Nas Organizações Acadêmicas, a proporção de presentes à prova que consideraram o Componente de Conhecimento Específico de grau de dificuldade Médio esteve entre 55,9 \% nos Centros Universitários e 63,9 \% nas Universidades ( ver também Tabela II.4, no Anexo II ).
O percentual de estudantes que considerou o Componente Específico da prova Difícil ou Muito difícil foi maior entre os estudantes de cursos de IES Privadas ( 39,3 \% ), uma diferença estatisticamente significativa da proporção em IES Públicas ( 27,7 \% ).
Já a proporção de presentes à prova que consideraram este componente de grau de dificuldade Médio foi 64,6 \% para os de IES Públicas e 56,3 \% para os de IES Privadas ( ver Gráfico 4.8 e Tabela II.4, no Anexo II ). 4.2 EXTENSÃO DA PROVA EM RELAÇÃO AO TEMPO TOTAL Indagados quanto à extensão da prova, em relação ao tempo total oferecido para a sua resolução ( Questão 3 ), os estudantes apontaram, com maior incidência, a alternativa que considerava a extensão adequada, para todas as agregações consideradas ( Gráfico 4.9, Gráfico 4.10 e, no Anexo II, a Tabela II.5 ).
O percentual de alunos que responderam ser a extensão da prova adequada foi de 67,9 \%.
Já 27,4 \% dos inscritos presentes consideraram que a prova foi Longa ou Muito longa, e 4,7 \% a avaliaram como Curta ou Muito curta.
Dentre as Grandes Regiões, a proporção daqueles que avaliaram a prova como Longa ou Muito longa, em relação ao tempo total destinado à sua resolução, variou de 23,7 \% na região Norte até 28,7 \% na região Sul.
Há diferença estatisticamente significativa apenas entre a região Norte e as regiões Sudeste e Sul.
Considerando -se o desempenho dos alunos, nota -se ainda que 64,1 \% consideraram a extensão da prova adequada, no quarto de desempenho inferior e 72,2 \%, no de melhor desempenho ( quarto superior ).
Nos quartos intermediários, essas proporções foram de 66,1 \%, no segundo quarto, e de 69,2 \%, no terceiro quarto, proporções crescentes com o desempenho.
No Gráfico 4.10, pode -se constatar que há diferença estatisticamente significativa das proporções de estudantes que consideraram a prova Longa ou Muito longa entre os dois primeiros quartos de desempenho e os dois últimos quartos, sendo que a diferença entre o primeiro e o segundo quartos e entre o terceiro e o quarto quartos não é estatisticamente significativa.
A maior proporção dessa resposta se encontra no primeiro quarto de desempenho ( 30,7 \% ), e a menor, no quarto superior ( 23,6 \% ), mas não caracterizando uma tendência.
Considerando -se a extensão da prova, o percentual de estudantes que avaliou a prova como Longa ou Muito longa, em relação ao tempo total destinado à sua resolução foi maior entre os estudantes de Centros Universitários, atingindo 30,0 \%, enquanto os de Universidades atingiram a menor proporção entre as Organização Acadêmica : 26,1 \%, sendo a diferença entre esses dois extremos a única que é estatisticamente significativa.
Nas Organizações Acadêmicas, a proporção de presentes à prova que considerou a sua extensão como adequada, em relação ao tempo total destinado à sua resolução esteve entre 64,5 \% nos Centros Universitários e 69,4 \% nas Universidades ( ver também Tabela II.6, no Anexo II ).
O percentual de estudantes que considerou a prova Longa ou Muito longa, em relação ao tempo total destinado à sua resolução foi maior entre os estudantes de cursos de IES Privadas ( 30,3 \% ), uma diferença estatisticamente significativa da proporção registrada em IES Públicas ( 25,5 \% ).
Já a proporção de presentes à prova que considerou a extensão Adequada foi de 69,8 \% para os de IES Públicas e de 65,1 \% para os de IES Privadas ( ver Gráfico 4.12 e Tabela II.6, no Anexo II ). 4.3 COMPREENSÃO DOS ENUNCIADOS DAS QUESTÕES 4.3.1 Componente de Formação Geral Em relação aos enunciados das questões do Componente de Formação Geral ( Questão 4 ), as opiniões foram positivas, já que 71,6 \% dos alunos avaliados consideraram todos ou a maioria dos enunciados das questões claros e objetivos ( Gráfico 4.13, Gráfico 4.14 e, no Anexo II, a Tabela II.7 ).
Na análise regional, a percentagem de estudantes que avaliou que todos ou a maioria dos enunciados das questões do Componente de Formação Geral estavam claros e objetivos variou de 69,9 \%, na região Sul, a 72,7 \%, na região Sudeste, sendo a diferença entre as proporções dessas regiões a única estatisticamente significativa.
A análise das percepções dos estudantes sobre a clareza e a objetividade dos enunciados permite afirmar que todos, ou a maioria dos enunciados de questões relativas ao Componente de Formação Geral, foram considerados claros e objetivos para a maior parte dos respondentes ( maior ou igual a 69,9 \% em todas as regiões e maior ou igual a 69,2 \% para todos os quartos de desempenho ).
Segundo o desempenho, observa -se que a proporção dos estudantes que emitiu essa opinião cresce conforme o desempenho aumenta, com diferença estatisticamente significativa apenas entre os quartos de desempenho extremos.
No quarto superior, a clareza e a objetividade de todos ou da maioria dos enunciados das questões foram percebidas por 73,6 \%.
A proporção dos presentes que considerou que todos ou a maioria dos enunciados das questões do Componente de Formação Geral estavam claros e objetivos foi maior entre os estudantes de Faculdades, atingindo 72,7 \%.
Já os de CEFET/IF alcançaram o menor índice : 69,4 \%.
No Gráfico 4.15, é possível observar que as diferenças entre todas as Organizações Acadêmicas não são estatisticamente significativas ( ver também Tabela II.8, no Anexo II ).
O percentual de estudantes que considerou que todos ou a maioria dos enunciados das questões do Componente de Formação Geral estavam claros e objetivos foi maior entre os estudantes de cursos de IES Privadas ( 71,8 \% ), sem diferença estatisticamente significativa da proporção registrada em IES Públicas ( 71,4 \% ).
Um maior detalhamento pode ser encontrado no Gráfico 4.16 e na Tabela II.8, no Anexo II. 4.3.2 Componente de Conhecimento Específico Em relação aos enunciados das questões do Componente de Conhecimento Específico da prova, 83,2 \% dos estudantes avaliados da Área de Agronomia, a clareza e a objetividade ( Questão 5 ) estavam presentes em todas ou na maioria das questões ( Gráfico 4.17 a Gráfico 4.20, e no Anexo II, as Tabelas II.9 e II.10 ).
A maioria dos estudantes de todas as Grandes Regiões brasileiras considerou claros e objetivos todos ou a maioria dos enunciados das questões do Componente de Conhecimento Específico da prova, percentual sempre maior ou igual a 82,7 \%.
As diferenças entre todas as regiões não são estatisticamente significativas.
A proporção de estudantes que considerou os enunciados das questões claros e objetivos apresenta uma tendência crescente em relação ao aumento de desempenho : mais elevada no quarto superior ( 87,8 \% ), se comparada ao quarto inferior de desempenho ( 77,1 \% ).
As diferenças entre todos os quartos de desempenho são estatisticamente significativas.
A proporção dos presentes que considerou que todos ou a maioria dos enunciados das questões do Componente de Conhecimento Específico estavam claros e objetivos foi maior entre os estudantes de Universidades, atingindo 84,3 \%.
Já os de Faculdades atingiram a menor proporção entre as Organizações Acadêmicas : 80,0 \%, sendo a diferença entre essas duas Organizações Acadêmicas a única que é estatisticamente significativa, como é possível observar no Gráfico 4.19 ( ver também Tabela II.10, no Anexo II ).
O percentual de estudantes que considerou que todos ou a maioria dos enunciados das questões do Componente de Conhecimento Específico estavam claros e objetivos foi maior entre os estudantes de cursos de IES Públicas ( 84,2 \% ), uma diferença estatisticamente significativa da proporção registrada em IES Privadas ( 81,6 \% ).
Ver Gráfico 4.20 e Tabela II.10, no Anexo II, para um maior detalhamento. 4.4 SUFICIÊNCIA DAS INFORMAÇÕES/INSTRUÇÕES FORNECIDAS Ao avaliarem as informações/instruções fornecidas para a resolução das questões ( Questão 6 ), 88,2 \% dos respondentes da Área de Agronomia de todo o Brasil afirmaram que estas eram até excessivas ou suficientes em todas ou na maioria das questões ( Gráfico 4.21, Gráfico 4.22, e, no Anexo II, a Tabela II.11 ).
Quanto à distribuição de respondentes pelas Grandes Regiões, observa -se que a proporção de estudantes que considerou as informações/instruções fornecidas até excessivas ou suficientes em todas ou na maioria das questões foi sempre superior ou igual a 85,0 \%, chegando a 89,4 \% na região Sul.
Existem diferenças estatisticamente significativas entre a região Norte e as regiões Sudeste, Sul e Centro-Oeste, e entre a região Nordeste e as regiões Sudeste, Sul e Centro-Oeste.
Levando -se em conta o desempenho dos participantes, observa -se uma tendência crescente das proporções de participantes que avaliaram as informações/instruções como até excessivas ou suficientes em todas ou na maioria das questões, havendo diferenças estatisticamente significativas entre todos os quartos de desempenho, como mostrado no Gráfico 4.22.
O percentual foi mais elevado no quarto superior ( 91,7 \% ), percentual superior à média nacional ( 88,2 \% ).
Já no quarto inferior, a suficiência das informações/instruções declarada como até excessiva, em todas ou na maioria das questões, foi percebida por 84,4 \% dos respondentes.
Considerando -se as informações/instruções fornecidas para a resolução das questões ( Questão 6 ), a proporção que considerou que estas eram até excessivas ou suficientes em todas ou na maioria das questões foi maior para os estudantes de Universidades, alcançando 88,9 \%, enquanto os de Faculdades registraram o menor índice : 86,5 \%, sendo a diferença entre essas duas Organizações Acadêmicas a única estatisticamente significativa, como é possível observar no Gráfico 4.23 ( ver também Tabela II.12, no Anexo II ).
O percentual de estudantes que considerou que as informações/instruções fornecidas para a resolução das questões eram até excessivas ou suficientes em todas ou na maioria das questões foi maior entre os estudantes de cursos de IES Públicas ( 88,9 \% ), uma diferença estatisticamente significativa da proporção registrada em IES Privadas ( 87,2 \% ).
Ver Gráfico 4.24 e Tabela II.12, no Anexo II, para maior detalhamento. 4.5 DIFICULDADE ENCONTRADA AO RESPONDER À PROVA Indagados sobre as dificuldades com as quais se depararam ao responder à prova ( Questão 7 ), 14,3 \% dos estudantes apontaram o Desconhecimento do conteúdo.
Para 52,0 \%, a Forma diferente de abordagem do conteúdo foi indicada como dificuldade.
Já a Falta de motivação para fazer a prova foi a dificuldade apontada por 8,7 \% dos respondentes.
Considerando -se todo o Brasil, 21,2 \% dos respondentes afirmaram que não tiveram qualquer tipo de dificuldade para responder à prova ( Tabelas II.13 e II.14 no Anexo II ).
Nos Gráficos 4.25 a 4.28, são apresentados os percentuais de estudantes que apontaram o Desconhecimento do conteúdo como dificuldade ao responder à prova.
Na análise por Grandes Regiões, os percentuais de inscritos e presentes que apontaram o Desconhecimento do conteúdo como dificuldade ao responder à prova não superaram 17,2 \%.
Os percentuais variaram de 13,7 \% na região Sul a 17,2 \%, na região Norte, sendo a diferença entre essa duas Grandes Regiões a única que foi estatisticamente significativa.
A Forma diferente de abordagem do conteúdo foi a escolha modal dos estudantes, com percentuais que variaram de 50,9 \% ( região Sul ) a 53,3 \% ( região Norte ).
Os percentuais de alunos que citaram a Falta de motivação como dificuldade variaram de 7,6 \% ( região Sul ) a 10,1 \% ( região Nordeste ).
Os que declararam não ter qualquer tipo de dificuldade para responder à prova variaram de 16,1 \% na região Norte a 23,7 \% na região Sul.
Em relação aos quartos de desempenho, o Desconhecimento do conteúdo foi a opção escolhida por 14,6 \% dos estudantes do quarto inferior e por 14,7 \% do quarto superior, não sendo observadas diferenças estatisticamente significativas entre os quartos de desempenho.
A alternativa modal para os alunos, quando agregados pelos quartos de desempenho, para a dificuldade encontrada, foi causada pela Forma diferente de abordagem do conteúdo, variando de 49,1 \% no quarto superior a 54,9 \% no segundo quarto.
Na análise por Organização Acadêmica, os percentuais de inscritos e presentes que apontaram o Desconhecimento do conteúdo como dificuldade ao responder à prova não superaram 14,9 \%.
Os percentuais variaram de 12,1 \% para os de CEFET/IF a 14,9 \%, para os de Universidades.
A Forma diferente de abordagem do conteúdo foi a escolha modal dos estudantes, com percentuais que variaram de 49,2 \% ( Universidades ) a 57,1 \% ( Faculdades ).
O percentual de alunos citando a Falta de motivação como dificuldade variou de 5,2 \% ( Centros Universitários ) a 10,6 \% ( Universidades ).
Os que declararam não ter qualquer tipo de dificuldade para responder à prova variaram de 18,6 \% nos CEFET/IF a 22,4 \% nos Centros Universitários.
No Gráfico 4.27, é possível observar que apenas a diferença entre Universidades e CEFET/IF é estatisticamente significativa ( ver também Tabela II.14, no Anexo II ).
Em relação às Categorias Administrativas, o Desconhecimento do conteúdo foi a opção escolhida por 14,5 \% dos estudantes nas IES Públicas e por 14,1 \% nas IES Privadas, sem diferença estatisticamente significativas entre essas duas categorias.
A alternativa modal para os alunos, quando agregados pelas Categorias Administrativas, para a dificuldade encontrada, foi causada pela Forma diferente de abordagem do conteúdo : 49,8 \% para os estudantes das \textbf{Públicas} e 55,5 \% para os das Privadas.
Ver Gráfico 4.28 e Tabela II.14, no Anexo II para maior detalhamento. 4.6 CONTEÚDOS DAS QUESTÕES OBJETIVAS DA PROVA Ao se analisarem os conteúdos das questões objetivas da prova ( Questão 8 ), constatou -se que um percentual muito pequeno dos estudantes avaliados, apenas 1,8 \%, afirmou que não estudou ainda a maioria desses conteúdos ( Gráficos 4.29 a Gráfico 4.32, e as Tabelas II.15 e II.16 no Anexo II ).
A maioria ( 80,5 \% ) afirmou ter estudado e aprendido muitos ou todos os conteúdos.
Na análise por Grande Região, a proporção de respondentes que escolheu a opção não estudou ainda a maioria desses conteúdos foi pequena.
Observa -se que, nas regiões Norte ( 2,5 \% ), Nordeste ( 2,7 \% ) e Centro-Oeste ( 2,3 \% ) as proporções foram maiores do que a média nacional ( 1,8 \% ).
Nas regiões Sudeste e Sul, a situação foi inversa : uma proporção mais baixa, 1,7 \% e 1,1 \% respectivamente.
São observadas diferenças estatisticamente significativas entre a região Sul e as regiões Norte, Nordeste e Centro-Oeste.
Em todas as regiões, a maioria dos presentes afirmou ter estudado e aprendido muitos ou todos os conteúdos, com proporções variando entre 71,5 \% na região Norte a 83,1 \% na região Sudeste.
Considerando -se separadamente as opiniões de estudantes dos quatro quartos de desempenho, observa -se que, no quarto inferior, 3,4 \% ofereceram como resposta que não estudou ainda a maioria desses conteúdos, sendo 0,6 \% os do quarto superior com a mesma resposta.
As proporções são decrescentes com o desempenho.
A diferença entre os alunos que optaram por esse motivo de dificuldade no quarto inferior em relação aos demais quartos é estatisticamente significativa, assim como entre o segundo e o último quartos.
Tendo -se em conta o quarto superior, 90,3 \% dos alunos afirmaram ter estudado e aprendido muitos ou todos os conteúdos.
No outro extremo, no primeiro quarto, 67,8 \% optaram pelas mesmas categorias.
As proporções são crescentes com o desempenho.
Na análise por Organização Acadêmica, as proporções de respondentes que escolheram a opção não estudou ainda a maioria desses conteúdos foram pequenas.
Observa -se que nos Centros Universitários ( 2,0 \% ), nas Faculdades ( 2,7 \% ) e nos CEFET/IF ( 2,1 \% ), as proporções foram maiores do que a média nacional ( 1,8 \% ).
Nas Universidades a situação foi inversa : registrou -se uma proporção mais baixa ( 1,5 \% ).
No Gráfico 4.31, é possível observar que apenas a diferença entre as Universidades e as Faculdades é estatisticamente significativa ( ver também Tabela II.16, no Anexo II ).
Em todos os tipos de Organização Acadêmica, a maioria dos presentes afirmou ter estudado e aprendido muitos ou todos os conteúdos, com proporções variando entre 76,2 \% nas Faculdades a 82,4 \% nas Universidades.
Em relação às Categorias Administrativas, observa -se que, nas IES Públicas, 2,0 \% ofereceram como resposta que não estudou ainda a maioria desses conteúdos, sendo 1,7 \% nas IES Privadas com a mesma resposta.
A diferença entre os alunos que optaram por esse motivo de dificuldade nas duas categorias não é estatisticamente significativa.
Considerando -se as IES Públicas, 81,7 \% dos alunos afirmaram ter estudado e aprendido muitos ou todos os conteúdos.
Na outra categoria, 78,7 \% optaram pelas mesmas categorias ( ver Gráfico 4.32 e Tabela II.16, no Anexo II ). 4.7 TEMPO GASTO PARA CONCLUIR A PROVA Ao responderem sobre o tempo de conclusão da prova ( Questão 9 ), mais de quatro quintos dos estudantes ( 81,3 \% ) afirmaram ter gasto Entre duas e quatro horas ( Gráfico 4.33 a Gráfico 4.36 e, no Anexo II, as Tabelas II.17 e II.18 ).
Considerando -se as cinco Grandes Regiões brasileiras, os percentuais dos que utilizaram Entre duas e quatro horas para finalizar a prova nas regiões Norte ( 85,8 \% ), Nordeste ( 84,0 \% ) e Sudeste ( 82,1 \% ) foram superiores ao percentual nacional.
Nas demais Grandes Regiões, o percentual de alunos que dispensaram Entre duas e quatro horas para concluir a prova ficou abaixo do percentual nacional : 78,5 \% na região Sul e 80,7 \% na região Centro-Oeste, como mostrado no Gráfico 4.33.
A diferença entre a região Norte e as regiões Sudeste, Sul e Centro-Oeste é estatisticamente significativa, assim como entre a região Nordeste e a região Sul, e entre as regiões Sudeste e Sul.
Analisadas as alternativas escolhidas pelos estudantes que se situam nos diferentes quartos de desempenho, observa -se uma tendência crescente : uma proporção maior de participantes no quarto superior ( 86,2 \% ) declarou ter gasto Entre duas e quatro horas para concluir a prova quando comparada com a proporção de participantes nos quartos inferiores.
Há diferenças estatisticamente significativas entre todos os quartos de desempenho.
Na análise por Organização Acadêmica, os percentuais dos que utilizaram Entre duas e quatro horas para finalizar a prova nas Universidades ( 79,7 \% ) foi inferior ao percentual nacional.
Nos demais tipos de Organização Acadêmica, o percentual de alunos que dispensaram Entre duas e quatro horas para concluir a prova ficou muito próximo : 84,6 nos Centros Universitários, 82,1 \% nas Faculdades e 83,8 \% nos CEFET/IF, como mostrado no Gráfico 4.35.
As diferenças entre as Universidades e os Centros Universitários e os CEFET/IF são estatisticamente significativas.
Em relação às Categorias Administrativas, observa -se que, nas IES Públicas, 80,2 \% declararam ter gasto Entre duas e quatro horas para concluir a prova, sendo 83,0 \% nas IES Privadas com a mesma resposta.
A diferença entre os alunos que optaram por esse motivo de dificuldade nas duas categorias é estatisticamente significativa ( ver Gráfico 4.36 e Tabela II.18, no Anexo II ).

% matched lemmas: pública
\end{textsample}
